\section{Limitations}
The primary limitation is evidence status: no independent reviewer has signed
off the RAGTruth or Diag-150 labels, and the human RQ4 study has not begun. All
headline metrics are therefore pre-review calibration evidence. The
LLM-simulated pilots use one model family under three prompts, not independent
humans or independent model populations. Their low RAGTruth agreement and poor
label-replacement sensitivity preclude using them as ground truth.

ContextTrace-Diag-150 is synthetic and designed around the implemented taxonomy.
The external RAGTruth experiment uses a 200-case deterministic sample rather
than the entire corpus. Answer-level labels do not always map cleanly to atomic
claim verdicts, which contributes to the 0.337 claim-verdict macro-F1. The
evidence-span metric is token overlap and does not directly measure boundary
minimality or multi-span recall.

The verifier uses deterministic lexical and semantic heuristics rather than a
general entailment model. It may fail on implicit reasoning, long dependency
chains, domain-specific numerical relations, multilingual traces, and evidence
distributed across many documents. Source trust and freshness require metadata
that many production traces omit. Grounding against supplied contexts does not
establish external factual truth.

NLI-only and judge-only ablations are unavailable because no qualifying frozen
model artifacts were registered. Latency measurements are workstation-specific.
The same-ID baseline cache does not expose every ContextTrace diagnostic and its
original generation cost is unavailable. Finally, the anonymous artifact can
reproduce software results but cannot redistribute human identities, private
condition keys, or future licensed response data.
